El saldo total administrado por el banco durante 2025 presentó fluctuaciones mensuales, sin evidenciar una tendencia sostenida de crecimiento o disminución a lo largo del período analizado.
La distribución de los saldos entre las diferentes sucursales fue relativamente equilibrada, sin observarse una concentración significativa en una única agencia. Las diferencias identificadas entre sucursales fueron moderadas en relación con el volumen total administrado.
De forma similar, los segmentos de clientes mostraron una distribución homogénea del saldo agregado, lo que sugiere que la cartera del banco se encuentra diversificada y no depende exclusivamente de un único perfil de clientes.
El análisis de dispersión no evidenció una relación lineal claramente definida entre el ingreso de los clientes y el saldo disponible en sus cuentas, lo que indica que el nivel de ingreso, por sí solo, no explica el comportamiento de los saldos observados.
La mayor concentración de saldos se registró en los clientes con edades comprendidas entre 20 y 69 años, mientras que los grupos de menor y mayor edad presentaron una participación considerablemente menor en el saldo total administrado.